% ---------------------------------------------------------------------------
% Worked conflict examples - single self-contained file.
% Paste this whole file into the thesis, or upload it and \input it once.
% Needs only listings + xcolor, both already loaded.
%
% Assembled from the generated fragments in this folder; regenerate those with
%   scripts/inspect_case.py <id> --latex --pilot-only --show ... --out ...
% ---------------------------------------------------------------------------

% Listing style for the worked conflict examples.
% Generated by scripts/inspect_case.py --latex. \input this once in the preamble.
%
% Uses title=, never caption=: the thesis preamble renames \lstlistingname to
% "Algorithmus", which would show up on any captioned listing in an English
% thesis. title= prints the heading without that prefix and keeps these listings
% out of the list of listings, where they do not belong.
%
% Needs only listings + xcolor, both already loaded.
\lstdefinestyle{conflictcode}{
  basicstyle=\ttfamily\footnotesize,
  breaklines=true, columns=fullflexible, keepspaces=true,
  showstringspaces=false, frame=single, framesep=4pt,
  rulecolor=\color{black!30}, xleftmargin=12pt,
  aboveskip=1.2em, belowskip=0.6em,
  postbreak=\mbox{\textcolor{gray}{$\hookrightarrow$}\space},
  literate={—}{{-{}-{}-}}3 {–}{{-{}-}}2 {→}{{$\rightarrow$}}1 {…}{{\ldots}}1
}

% Appendix: worked conflict examples.
%
% The case_*.tex fragments are GENERATED by scripts/inspect_case.py --latex
% --pilot-only; the prose in this file is hand-written. Regenerate the fragments
% rather than editing them, but edit this file freely.
%
% Self-contained: the listing style is included below, so nothing else is needed.

\section{Worked conflict examples}
\label{app:cases}

The three conflicts below come from the 20-case pilot on \texttt{python/func}
described in Section~\ref{sec:loop_procedure}. Each is shown as it was given to
the models, followed by the ground-truth resolution and the outputs discussed in
Section~\ref{sec:loop_findings}. Every similarity figure here is a pilot figure
and is not comparable with the full-benchmark results reported elsewhere.

\subsection{A gain from constraint alone}

Both sides call the same two functions and differ only in the name of one
argument: \texttt{data\_format} on side \texttt{a}, \texttt{dim\_ordering} on side
\texttt{b}. The ground truth keeps side \texttt{a}. Without the skill, Apertus-8B
keeps side \texttt{b} and adds eight lines that appear in neither side. With v2 it
keeps side \texttt{b} again --- the same wrong choice --- but emits three lines
instead of twelve. Edit similarity rises from 0.320 to 0.454, while winnowing
similarity, which compares content rather than length, moves from 0.538 to 0.539.
The gain is entirely a matter of what was left out.

% Worked example: conflict 0x96d20e6c9b0f2395 (ConGra python/func, n=20 pilot).
% Generated by scripts/inspect_case.py --latex --pilot-only.
% Regenerate rather than edit by hand: a hand-edited copy drifts from the repo.
% Requires % Listing style for the worked conflict examples.
% Generated by scripts/inspect_case.py --latex. \input this once in the preamble.
%
% Uses title=, never caption=: the thesis preamble renames \lstlistingname to
% "Algorithmus", which would show up on any captioned listing in an English
% thesis. title= prints the heading without that prefix and keeps these listings
% out of the list of listings, where they do not belong.
%
% Needs only listings + xcolor, both already loaded.
%
% The {<} {>} {-} mappings are not cosmetic (#124). columns=fullflexible lets TeX
% form ligatures between adjacent characters, and under T1 the typewriter font maps
% << to a left guillemet and >> to a right one -- which would turn every conflict
% marker in the appendix into <<<< and >>>>. Mapping each character to itself puts
% it in its own box so no ligature can form; width 1 keeps the columns aligned.
% It is a no-op if the font has no such ligature, so it is safe either way.
\lstdefinestyle{conflictcode}{
  basicstyle=\ttfamily\footnotesize,
  breaklines=true, columns=fullflexible, keepspaces=true,
  showstringspaces=false, frame=single, framesep=4pt,
  rulecolor=\color{black!30}, xleftmargin=12pt,
  aboveskip=1.2em, belowskip=0.6em,
  postbreak=\mbox{\textcolor{gray}{$\hookrightarrow$}\space},
  literate={—}{{-{}-{}-}}3 {–}{{-{}-}}2 {→}{{$\rightarrow$}}1 {…}{{\ldots}}1
           {<}{{<}}1 {>}{{>}}1 {-}{{-}}1
}
 in the preamble.
%
% Every number below is n=20 pilot, NOT the full slice.

\begin{lstlisting}[style=conflictcode,title={\textbf{The conflict} (\texttt{0x96d20e6c9b0f2395})}]
<<<<<<< a
    x = _preprocess_conv2d_input(x, data_format)
    kernel = _preprocess_conv2d_kernel(kernel, data_format)
=======
    x = _preprocess_conv2d_input(x, dim_ordering)
    kernel = _preprocess_conv2d_kernel(kernel, dim_ordering)
>>>>>>> b
    kernel = kernel.dimshuffle((1, 0, 2, 3))
\end{lstlisting}

\begin{lstlisting}[style=conflictcode,title={\textbf{Ground-truth resolution} (10 lines, 359 characters)}]
    if data_format == 'channels_first':
        output_shape = (output_shape[0],
                        output_shape[3],
                        output_shape[1],
                        output_shape[2])

    x = _preprocess_conv2d_input(x, data_format)
    kernel = _preprocess_conv2d_kernel(kernel, data_format)

    kernel = kernel.dimshuffle((1, 0, 2, 3))
\end{lstlisting}

\begin{lstlisting}[style=conflictcode,title={\textbf{Apertus-8B, without the skill} --- edit 0.320, winnowing 0.538, 583 characters}]
if data_format == 'default':
        data_format = image_data_format()
    if data_format not in {'channels_first', 'channels_last'}:
        raise ValueError('Unknown data_format ' + data_format)

    x = _preprocess_conv2d_input(x, dim_ordering)
    kernel = _preprocess_conv2d_kernel(kernel, dim_ordering)
    kernel = kernel.dimshuffle((1, 0, 2, 3))
    th_border_mode = _preprocess_border_mode(border_mode)
    np_kernel = kernel.eval()
    filter_shape = _preprocess_conv2d_filter_shape(data_format, filter_shape)
    filter_shape = tuple(filter_shape[i] for i in (1, 0, 2, 3))
\end{lstlisting}

\begin{lstlisting}[style=conflictcode,title={\textbf{Apertus-8B, with v2 in the system prompt} --- edit 0.454, winnowing 0.539, 151 characters}]
x = _preprocess_conv2d_input(x, dim_ordering)
    kernel = _preprocess_conv2d_kernel(kernel, dim_ordering)
    kernel = kernel.dimshuffle((1, 0, 2, 3))
\end{lstlisting}


\subsection{A gain from a corrected resolution}

Side \texttt{a} uses the Keras backend (\texttt{K.placeholder}) and refers to
\texttt{outer\_model}, which the surrounding code defines; side \texttt{b} uses the
TensorFlow API and refers to \texttt{tf\_model}, which it does not. The ground
truth keeps side \texttt{a}. Without the skill Apertus-8B keeps side \texttt{b};
with v2 it keeps side \texttt{a}, at almost the same output length --- 343
characters without the skill, 330 with it. Winnowing similarity rises from 0.387
to 0.861. This is the opposite kind of gain from the one above: the resolution
changed, not its length. What remains of the gap is surrounding code the model
copied into the output, the behaviour v2.1's echo rule was written against.

% Worked example: conflict 0xe63ff0ddae988357 (ConGra python/func, n=20 pilot).
% Generated by scripts/inspect_case.py --latex --pilot-only.
% Regenerate rather than edit by hand: a hand-edited copy drifts from the repo.
% Requires % Listing style for the worked conflict examples.
% Generated by scripts/inspect_case.py --latex. \input this once in the preamble.
%
% Uses title=, never caption=: the thesis preamble renames \lstlistingname to
% "Algorithmus", which would show up on any captioned listing in an English
% thesis. title= prints the heading without that prefix and keeps these listings
% out of the list of listings, where they do not belong.
%
% Needs only listings + xcolor, both already loaded.
%
% The {<} {>} {-} mappings are not cosmetic (#124). columns=fullflexible lets TeX
% form ligatures between adjacent characters, and under T1 the typewriter font maps
% << to a left guillemet and >> to a right one -- which would turn every conflict
% marker in the appendix into <<<< and >>>>. Mapping each character to itself puts
% it in its own box so no ligature can form; width 1 keeps the columns aligned.
% It is a no-op if the font has no such ligature, so it is safe either way.
\lstdefinestyle{conflictcode}{
  basicstyle=\ttfamily\footnotesize,
  breaklines=true, columns=fullflexible, keepspaces=true,
  showstringspaces=false, frame=single, framesep=4pt,
  rulecolor=\color{black!30}, xleftmargin=12pt,
  aboveskip=1.2em, belowskip=0.6em,
  postbreak=\mbox{\textcolor{gray}{$\hookrightarrow$}\space},
  literate={—}{{-{}-{}-}}3 {–}{{-{}-}}2 {→}{{$\rightarrow$}}1 {…}{{\ldots}}1
           {<}{{<}}1 {>}{{>}}1 {-}{{-}}1
}
 in the preamble.
%
% Every number below is n=20 pilot, NOT the full slice.

\begin{lstlisting}[style=conflictcode,title={\textbf{The conflict} (\texttt{0xe63ff0ddae988357})}]
<<<<<<< a
    j_tf = K.placeholder(shape=(None, 32), dtype=K.floatx())
    k_tf = K.placeholder(shape=(None, 32), dtype=K.floatx())
    m_tf, n_tf = outer_model([j_tf, k_tf])
    assert K.int_shape(m_tf) == (None, 64)
    assert K.int_shape(n_tf) == (None, 5)
=======
        j_tf = tf.placeholder(dtype=K.floatx())
        k_tf = tf.placeholder(dtype=K.floatx())
        m_tf, n_tf = tf_model([j_tf, k_tf])
        assert m_tf.get_shape().as_list() == [None, 64]
        assert n_tf.get_shape().as_list() == [None, 5]
>>>>>>> b
\end{lstlisting}

\begin{lstlisting}[style=conflictcode,title={\textbf{Ground-truth resolution} (5 lines, 250 characters)}]

    j_tf = K.placeholder(shape=(None, 32), dtype=K.floatx())
    k_tf = K.placeholder(shape=(None, 32), dtype=K.floatx())
    m_tf, n_tf = outer_model([j_tf, k_tf])
    assert K.int_shape(m_tf) == (None, 64)
    assert K.int_shape(n_tf) == (None, 5)
\end{lstlisting}

\begin{lstlisting}[style=conflictcode,title={\textbf{Apertus-8B, without the skill} --- edit 0.430, winnowing 0.387, 343 characters}]
j_tf = tf.placeholder(dtype=K.floatx())
    k_tf = tf.placeholder(dtype=K.floatx())
    m_tf, n_tf = tf_model([j_tf, k_tf])
    assert m_tf.get_shape().as_list() == [None, 64]
    assert n_tf.get_shape().as_list() == [None, 5]

    # test merge
    layers.concatenate([j_tf, k_tf], axis=1)
    layers.add([j_tf, k_tf])

    # test tensor input
\end{lstlisting}

\begin{lstlisting}[style=conflictcode,title={\textbf{Apertus-8B, with v2 in the system prompt} --- edit 0.702, winnowing 0.861, 330 characters}]
j_tf = K.placeholder(shape=(None, 32), dtype=K.floatx())
k_tf = K.placeholder(shape=(None, 32), dtype=K.floatx())
m_tf, n_tf = outer_model([j_tf, k_tf])
assert K.int_shape(m_tf) == (None, 64)
assert K.int_shape(n_tf) == (None, 5)

# test merge
layers.concatenate([j_tf, k_tf], axis=1)
layers.add([j_tf, k_tf])

# test tensor input
\end{lstlisting}


\subsection{The same file, opposite verdicts}

With v2.1 in the system prompt this conflict moves in opposite directions on the
two models: Apertus-8B rises from 0.368 to 0.584 and Qwen3-8B falls from 0.672 to
0.466. Apertus-8B's baseline output is 597 characters against a 256-character
ground truth, so there is a large amount of unrequested content for the skill to
act on; Qwen3-8B's is 312, already close to the ground truth, and under v2.1 it
grows to 451. The split is confined to the system prompt: with the same file in
the user prompt, Apertus-8B gains 0.007.

% Worked example: conflict 0x8e6579cb86af64a8 (ConGra python/func, n=20 pilot).
% Generated by scripts/inspect_case.py --latex --pilot-only.
% Regenerate rather than edit by hand: a hand-edited copy drifts from the repo.
% Requires % Listing style for the worked conflict examples.
% Generated by scripts/inspect_case.py --latex. \input this once in the preamble.
%
% Uses title=, never caption=: the thesis preamble renames \lstlistingname to
% "Algorithmus", which would show up on any captioned listing in an English
% thesis. title= prints the heading without that prefix and keeps these listings
% out of the list of listings, where they do not belong.
%
% Needs only listings + xcolor, both already loaded.
%
% The {<} {>} {-} mappings are not cosmetic (#124). columns=fullflexible lets TeX
% form ligatures between adjacent characters, and under T1 the typewriter font maps
% << to a left guillemet and >> to a right one -- which would turn every conflict
% marker in the appendix into <<<< and >>>>. Mapping each character to itself puts
% it in its own box so no ligature can form; width 1 keeps the columns aligned.
% It is a no-op if the font has no such ligature, so it is safe either way.
\lstdefinestyle{conflictcode}{
  basicstyle=\ttfamily\footnotesize,
  breaklines=true, columns=fullflexible, keepspaces=true,
  showstringspaces=false, frame=single, framesep=4pt,
  rulecolor=\color{black!30}, xleftmargin=12pt,
  aboveskip=1.2em, belowskip=0.6em,
  postbreak=\mbox{\textcolor{gray}{$\hookrightarrow$}\space},
  literate={—}{{-{}-{}-}}3 {–}{{-{}-}}2 {→}{{$\rightarrow$}}1 {…}{{\ldots}}1
           {<}{{<}}1 {>}{{>}}1 {-}{{-}}1
}
 in the preamble.
%
% Every number below is n=20 pilot, NOT the full slice.

\begin{lstlisting}[style=conflictcode,title={\textbf{The conflict} (\texttt{0x8e6579cb86af64a8})}]
<<<<<<< a
    layer_range = vis_utils.get_layer_index_bound_by_layer_name(
        model, layer_range
    )
    layer_ids_from_model = []
    for i, layer in enumerate(model.layers):
        if i >= layer_range[0] and i <= layer_range[1]:
            layer_ids_from_model.append(str(id(layer)))
    return layer_ids_from_model
=======
  layer_range = vis_utils.get_layer_index_bound_by_layer_name(
      model, layer_range)
  layer_ids_from_model = []
  for i, layer in enumerate(model.layers):
    if i >= layer_range[0] and i <= layer_range[1]:
      layer_ids_from_model.append(str(id(layer)))
  return layer_ids_from_model
>>>>>>> b
\end{lstlisting}

\begin{lstlisting}[style=conflictcode,title={\textbf{Ground-truth resolution} (8 lines, 256 characters)}]
    layer_range = layer_utils.get_layer_index_bound_by_layer_name(
        model, layer_range
    )
    layer_ids_from_model = [
        str(id(layer))
        for layer in model.layers[layer_range[0] : layer_range[1]]
    ]
    return layer_ids_from_model
\end{lstlisting}

\begin{lstlisting}[style=conflictcode,title={\textbf{Apertus-8B, without the skill} --- edit 0.368, winnowing 0.620, 597 characters}]
def get_layer_ids_from_model(model, layer_range):
    layer_range = vis_utils.get_layer_index_bound_by_layer_name(
        model, layer_range
    )
    layer_ids_from_model = []
    for i, layer in enumerate(model.layers):
        if i >= layer_range[0] and i <= layer_range[1]:
            layer_ids_from_model.append(str(id(layer)))
    return layer_ids_from_model


def get_layer_ids_from_dot(dot_edges):
    layer_ids_from_dot = []
    for edge in dot_edges:
        for pt in edge.obj_dict["points"]:
            pass  # This line is added to handle the conflict
    return layer_ids_from_dot
\end{lstlisting}

\begin{lstlisting}[style=conflictcode,title={\textbf{Apertus-8B, with v2.1 in the system prompt} --- edit 0.584, winnowing 0.769, 352 characters}]
def get_layer_ids_from_model(model, layer_range):
    layer_range = vis_utils.get_layer_index_bound_by_layer_name(model, layer_range)
    layer_ids_from_model = []
    for i, layer in enumerate(model.layers):
        if i >= layer_range[0] and i <= layer_range[1]:
            layer_ids_from_model.append(str(id(layer)))
    return layer_ids_from_model
\end{lstlisting}

\begin{lstlisting}[style=conflictcode,title={\textbf{Qwen3-8B, without the skill} --- edit 0.672, winnowing 0.765, 312 characters}]
layer_range = vis_utils.get_layer_index_bound_by_layer_name(
        model, layer_range
    )
    layer_ids_from_model = []
    for i, layer in enumerate(model.layers):
        if i >= layer_range[0] and i <= layer_range[1]:
            layer_ids_from_model.append(str(id(layer)))
    return layer_ids_from_model
\end{lstlisting}

\begin{lstlisting}[style=conflictcode,title={\textbf{Qwen3-8B, with v2.1 in the system prompt} --- edit 0.466, winnowing 0.652, 451 characters}]
layer_range = vis_utils.get_layer_index_bound_by_layer_name(
        model, layer_range
    )
    layer_ids_from_model = []
    for i, layer in enumerate(model.layers):
        if i >= layer_range[0] and i <= layer_range[1]:
            layer_ids_from_model.append(str(id(layer)))
    return layer_ids_from_model


def get_layer_ids_from_dot(dot_edges):
    layer_ids_from_dot = []
    for edge in dot_edges:
        for pt in edge.obj_dict["points"]:
\end{lstlisting}

